\begin{abstract}
Weak surface exchange can produce large variations in the spreading of
tracers among nominally similar channels. We study how a fixed integral
budget of tangential surface diffusivity should be distributed when the
exchange pattern is uncertain, and how accurately that pattern must be
measured before placement. For the dimensionless family
$k_c(s)=(c+\cos s)^2$, with $c$ uniform on $[-2,2]$, we prove sharp
small-budget asymptotics over all nonnegative integrable mobility fields.
Predetermined spatial design moves the disorder-moment threshold from
$q=4/3$ for uniform mobility to $q=8/5$, where the optimal moment is
$K_{\rm crit}M^{-2/5}[\log(1/M)]^{7/5}$. Above this threshold an
attained whole-line variational problem supplies the sharp coefficient.
The three optimized orders persist for smooth families with transverse
folds. Exact observation changes the optimal mean from order $M^{-1/4}$
to $M^{-1/5}$. For equal observation bins of width $\Delta$, we prove
the joint law $\mathcal Q\asymp M^{-1/5}+M^{-1/4}\Delta^{1/4}$,
an exact continuous crossover at $\Delta/M^{1/5}=O(1)$, and sharp
endpoint coefficients. A bulk--surface comparison transfers these
scalar optimum laws to stationary flow-induced dispersion at fixed
positive bulk diffusivity. Explicit local certificates, unrestricted
lower bounds, and numerical refinement distinguish asymptotic values
from the performance of finite-budget trial designs.
\end{abstract}

\section{Introduction}
\label{sec:introduction}
Particles transported through an adsorbing channel alternate between
advection in the bulk and periods of residence at the wall. Even when
the mean fraction of adsorbed particles is fixed, slow exchange at some
wall locations can increase fluctuations in the time spent moving.
Surface diffusion provides an escape route along the wall. The design
question is therefore spatial: where should a limited amount of surface
mobility be placed when the locations and strengths of slow exchange
are uncertain?

Adsorption, surface motion, and their effects on dispersion are established
parts of generalized Taylor dispersion theory
\citep{DillBrenner1982,LevesqueEtAl2012}. More recent formulations allow
spatially varying diffusivity and wall interactions in confined
microfluidic systems \citep{AlexandreGuerinDean2021}. The present problem
keeps the channel geometry, bulk flow, and equilibrium adsorption affinity
fixed. The desorption rate varies along the wall, while adsorption is
proportional to it with fixed proportionality. Consequently the mean
speed remains fixed, and the design changes the spreading about that
mean. The uncertainty is quenched: one kinetic pattern is selected for
each channel and stays fixed during transport. Moments over that
uncertainty describe reproducibility among channels.

Our model family on a wall of dimensionless length $2\pi$ is
\[
       k_c(s)=(c+\cos s)^2,\qquad c\sim\mathrm{Uniform}[-2,2].
\]
For $|c|<1$ the rate has two quadratic zeros. They merge into a
quartic zero at $c=\pm1$ and disappear for $|c|>1$.
We refer to this merger as a \emph{fold}. The family is deliberately
simple enough to retain exact constants, while containing both ordinary
kinetic defects and rare nearly merged defects. A designer chooses a
nonnegative tangential diffusivity $D(s)$ with
$\int_\Gamma D\dd s=M$. The budget is dimensionless in fixed physical
units specified in \cref{sec:model}; no powers of a dimensional budget
are intended. No mobility cap or minimum fabrication length is imposed.

\subsection{Results and the role of information}
For a fixed pattern and mobility, the singular wall response is the
variational quantity
\begin{equation}\label{eq:intro-response}
 J_c(D)=\sup_{v\in C^\infty(\Gamma)}
 \left\{2\int_\Gamma v\dd s
      -\int_\Gamma[D|v'|^2+k_cv^2]\dd s\right\}.
\end{equation}
When an inverse exists this equals the integral of the solution of
$-\partial_s(D\partial_sh)+k_ch=1$. Its variational definition also
allows arbitrary integrable coefficients with zero sets or fine structure.
For physically admissible designs with finite response, the
flow-induced dispersion equals $\chi J_c$ plus a nonnegative bulk
correction, where $\chi=KV^2/(A+KP)$ is fixed by equilibrium and
mean flow. We prove a logarithmic bound on that correction after adding
an arbitrarily small uniform part to the mobility. This construction
preserves both the budget and the information used to choose the field.
It makes the leading scalar optimal values physically attainable,
without taking an infinite-bulk-mixing limit.

There are two related design questions. First, a predetermined field
minimizes the positive moment $\E[J_c(D)^q]$ before $c$ is measured.
Second, a field selected after observation minimizes the mean, with
exactly budget $M$ for every observation. The resulting moment orders
are compared in \cref{tab:regimes}; all their scalar coefficients are
given by sharp equivalents in \cref{thm:uniform-moments,thm:blind-sharp}.
In particular, the supercritical coefficient is an attained variational
value, rather than the cost of one chosen profile.

\begin{table}[t]
\centering
\caption{Small-budget orders of unrooted disorder moments in the cosine
ensemble. Uniform mobility has $D=M/(2\pi)$; predetermined design is
optimized over every nonnegative integrable field of budget $M$.
The threshold column separates the two open ranges in each row.
All orders have sharp coefficients in the cited theorems.}
\label{tab:regimes}
\renewcommand{\arraystretch}{1.25}
\begin{tabular}{@{}lclll@{}}
\toprule
Design & Threshold $q_*$ & $q<q_*$ & $q=q_*$ & $q>q_*$\\
\midrule
Uniform & $4/3$ & $M^{-q/4}$ & $M^{-1/3}\log(1/M)$
 & $M^{1/3-q/2}$\\
Predetermined & $8/5$ & $M^{-q/4}$
 & $M^{-2/5}[\log(1/M)]^{7/5}$ & $M^{(2-3q)/7}$\\
\bottomrule
\end{tabular}
\end{table}

The change of threshold has a direct interpretation. For low moments,
the leading cost comes from ordinary roots moving over the wall. The
formal normalized allocation is proportional to
$|\sin s|^{-\alpha_q}$, where $\alpha_q=(6q-4)/(q+4)$.
At $q=8/5$ this shape ceases to be integrable. The actual design must
then account for budget-dependent neighborhoods of the known fold sites.
The proof makes this argument valid against arbitrary concentrated or
oscillatory competitors. It gives the exact critical logarithm and,
above the threshold, matching localization and recovery through a
whole-line fold problem. The same three orders hold for a class of
$C^4$ squared-rate families with isolated transverse folds
(\cref{thm:generic-orders}). The generic result is an order theorem;
it does not assign the cosine coefficients to other families.

Observation changes the mean more substantially than predetermined
grading alone. Exact knowledge of $c$ permits placement at the realized
zeros and gives
$\mathcal O(M)\sim K_{\rm obs}M^{-1/5}$, whereas
$\mathcal P_1(M)\sim K_1M^{-1/4}$.
For equal bins of width $\Delta=4/N$, let $\mathcal Q(M,\Delta)$
be the optimal mean when only the bin index is revealed.
\Cref{thm:precision} proves the uniform two-parameter law
\[
 \mathcal Q(M,\Delta)\asymp
          M^{-1/5}+M^{-1/4}\Delta^{1/4},\qquad M,\Delta\downarrow0.
\]
It also gives an exact crossover:
$M^{1/5}\mathcal Q(M,\Delta)\to\mathcal H(\tau)$ when
$\Delta/M^{1/5}\to\tau<\infty$. The continuous function
$\mathcal H$ is an explicit integral of the attained unit-budget
problem for a quadratic defect with an uncertain center.
It has $\mathcal H(0)=K_{\rm obs}$ and
$\mathcal H(\tau)\sim K_{\rm coarse}\tau^{1/4}$ at large $\tau$;
the simultaneous coarse limit is proved separately.
Thus offset resolution of order $M^{1/5}$ is the scale at which the
measurement changes the leading mean performance. This is a statement
about a specified noiseless quantizer, not a theory of arbitrary noisy
measurements or an optimization of the sensor.

\subsection{Relation to existing design and uncertainty methods}
The response \eqref{eq:intro-response} has the variational structure
of compliance. Mass-constrained reinforcement, measure relaxation,
and gradient constraints are established tools in this setting
\citep{BouchitteButtazzo2001,ButtazzoOudetVelichkov2015}.
Deterministic coefficient design under random forcing is treated by
\citet{ButtazzoMaestre2011}; general risk criteria for mixtures of
two positive conductors are studied by
\citet{AlphonseKunstekVrdoljak2026}. Those works provide close
optimization precedents. The present singular limit instead concerns
random vanishing reaction rates and an unrestricted nonnegative
$L^1$ diffusivity whose total budget tends to zero. The two cited
stochastic-conductivity models impose positive conductivity bounds
and do not include this coefficient class.
The contribution here is the sharp transport laws, including the
unrestricted lower bounds and recovery constructions at the degeneracies.

Neither the relation between risk preference and resource allocation
nor the importance of rare degeneracies is a new general principle.
Optimized tolerance models already show that risk preference changes
large-loss statistics \citep{NewmanGirvanFarmer2002}. Parameter-averaged
moments can likewise select rare bifurcations because their enhanced
response compensates for their small probability
\citep{BerryKeatingSchomerus2000}. Here the integral mobility budget
and local resolvent identify the particular thresholds, logarithm,
and attained variational constants.

The information available at the design stage is also an established
part of stochastic shape optimization \citep{ContiEtAl2009}.
Observation-channel continuity and quantized policies have a developed
theory \citep{YukselLinder2012,SaldiYukselLinder2017}.
The result proved here is a quantitative joint limit in which
observation resolution and an independent, degenerating PDE budget
both vary. Conditional localization must hold uniformly, and bins
containing merged defects must be controlled before the local values
can be averaged. These are the additional steps that produce
$\mathcal H$ and the global resolution threshold.

\subsection{Organization}
\Cref{sec:model} derives the physical coefficient, specifies rough
mobility realizations, and proves the scalar-to-bulk comparison.
\Cref{sec:local-baseline} establishes the harmonic and fold responses,
the uniform-design moments, and the limits of replacing random rates
by averages or local zero statistics. \Cref{sec:predetermined} proves
the sharp predetermined-design theorem; \cref{sec:generic-folds}
extends its orders to generic folds. \Cref{sec:exact-observation}
derives the exact local placement profile and the observation-dependent
mean. \Cref{sec:finite-precision} proves the finite-resolution theorem.
\Cref{sec:numerics} supplies reproducible numerical evidence with
separate discretization checks, and \cref{sec:discussion} discusses
the physical scope. All arguments needed for these results are included.
